% ============================================================
%  Equações para justificação do HLM vs OLS
%  TCC — Ricardo Calheiros — ESALQ/USP — 2026
% ============================================================

% ----------------------------------------------------------
%  Equação 1: Modelo OLS de referência (linha de base)
% ----------------------------------------------------------
\begin{equation}
  \ln(W)_{ijk} = \alpha_0
    + \beta_1 \cdot \text{Negro}_{ijk}
    + \mathbf{X}_{ijk}'\boldsymbol{\beta}
    + \varepsilon_{ijk}, \quad
  \varepsilon_{ijk} \overset{\text{iid}}{\sim} \mathcal{N}\!\left(0,\,\sigma^2\right)
  \label{eq:ols}
\end{equation}

\noindent
\textbf{Limitação do OLS:} assume $\mathrm{Cov}(\varepsilon_{ijk},\varepsilon_{ij'k})=0$ para
todo $j \neq j'$ dentro do mesmo estado $k$. Quando existe agrupamento (clustering), essa
suposição é violada: observações do mesmo estado compartilham um contexto comum não modelado,
tornando os erros-padrão do OLS inconsistentes (subestimados ou super estimados) mesmo com
correção por \textit{cluster}.

% ----------------------------------------------------------
%  Equação 2: Modelo Nulo HLM (M0) — decomposição de variância
% ----------------------------------------------------------
\begin{equation}
  \ln(W)_{ijk} = \delta_{000} + u_{00k} + \varepsilon_{ijk}, \quad
  u_{00k} \sim \mathcal{N}\!\left(0,\,\tau^2_v\right), \quad
  \varepsilon_{ijk} \sim \mathcal{N}\!\left(0,\,\sigma^2\right)
  \label{eq:m0}
\end{equation}

% ----------------------------------------------------------
%  Equação 3: Coeficiente de Correlação Intraclasse (ICC)
% ----------------------------------------------------------
\begin{equation}
  \rho_{UF} = \frac{\hat{\tau}^2_v}{\hat{\tau}^2_v + \hat{\sigma}^2}
  = \frac{0.08342}{0.08342 + 0.7653}
  = 0.0983
  \label{eq:icc_valor}
\end{equation}

\noindent
O valor $\rho_{UF} = 9.83\%$ supera o limiar de $5\%$ recomendado por
\citet{raudenbush2002}, confirmando que o nível~3 (UF) explica variância não negligenciável
do log-rendimento e justifica a estrutura multinível.

% ----------------------------------------------------------
%  Equação 4: Efeito de design (Design Effect)
% ----------------------------------------------------------
\begin{equation}
  \text{DEFF} = 1 + \left(\bar{n} - 1\right)\,\rho_{UF}
  \approx 1 + \left(56,959 - 1\right) \times 0.0983
  \approx 5,600
  \label{eq:deff}
\end{equation}

\noindent
O efeito de design $\approx 5,600$ indica que a variância amostral das estimativas OLS
é inflada por esse fator quando a estrutura de clustering é ignorada. Erros-padrão OLS sem
correção de cluster estariam subestimados por $\sqrt{\text{DEFF}} \approx 75\times$.

% ----------------------------------------------------------
%  Equação 5: Teste de Wald para H0: τ²_UF = 0 (M0)
% ----------------------------------------------------------
\begin{equation}
  z_{M0} = \frac{\hat{\tau}^2_v}{\widehat{\text{SE}}(\hat{\tau}^2_v)}
  = \frac{0.08342}{0.02314}
  = 3.61, \quad p_\text{boundary} = 0.0002
  \label{eq:wald_m0}
\end{equation}

\begin{equation}
  z_{M1} = \frac{\hat{\tau}^2_v}{\widehat{\text{SE}}(\hat{\tau}^2_v)}
  = \frac{0.06109}{0.01694}
  = 3.61, \quad p_\text{boundary} = 0.0002
  \label{eq:wald_m1}
\end{equation}

\noindent
O teste de fronteira (\textit{boundary test}) utiliza a distribuição mista
$\frac{1}{2}\chi^2_0 + \frac{1}{2}\chi^2_1$ (cf.\ \citealt{raudenbush2002}) adequada para
hipóteses nulas na fronteira do espaço paramétrico ($\tau^2 \geq 0$).
Em M0, $z = 3.61$ ($p < 0{,}001$***); em M1, $z = 3.61$
($p < 0{,}001$***): $\hat{\tau}^2_{UF}$ é estatisticamente diferente de zero
em todos os modelos.

% ----------------------------------------------------------
%  Equação 6: LRT entre modelos sequenciais (HLM REML)
% ----------------------------------------------------------
\begin{align}
  \Lambda_{M0 \to M1} &= -2\left(\ell_{M0} - \ell_{M1}\right)
    = -2\left(-1,976,646.50 - (-1,830,562.34)\right)
    = 292,168.3 \label{eq:lrt01} \\
  \Lambda_{M1 \to M2} &= -2\left(\ell_{M1} - \ell_{M2}\right)
    = 72,470.8 \label{eq:lrt12}
\end{align}

\noindent
Ambas as estatísticas superam em várias ordens de grandeza o valor crítico
$\chi^2_{(8,\,0{,}001)} = 26{,}1$ e $\chi^2_{(3,\,0{,}001)} = 16{,}3$,
confirmando que cada nível hierárquico adiciona poder explicativo significativo.

% ----------------------------------------------------------
%  Tabela resumo: HLM M0 — componentes de variância
% ----------------------------------------------------------
\begin{table}[H]
\centering
\caption{Componentes de variância e testes de significância — HLM de Três Níveis.
  $\hat{\sigma}^2$: variância Nível~1 (intraindividual).
  $\hat{\tau}^2_{UF}$: variância Nível~3 (inter-UF).
  ICC$_{UF}$: fração da variância total explicada pelo estado.
  Teste de Wald (boundary): H$_0$: $\tau^2 = 0$.}
\label{tab:variancia_componentes}
\begin{tabular}{lcccccc}
\toprule
\textbf{Modelo} & $\hat{\sigma}^2$ & $\hat{\tau}^2_{UF}$ & SE($\hat{\tau}^2_{UF}$) & $z$ & ICC$_{UF}$ & Sig. \\
\midrule
M0 (Nulo)         & 0.7653 & 0.08342 & 0.02314 & 3.61 & 9.83\% & *** \\
M1 (Individual)   & 0.6329 & 0.06109 & 0.01694 & 3.61 & 8.80\% & *** \\
M2 (+ UPA)        & 0.6037 & 0.03365 & 0.00933 & 3.61 & 5.28\% & *** \\
M3 (Completo)     & 0.6037 & 0.02504 & 0.00694 & 3.61 & 3.98\% & *** \\
\bottomrule
\multicolumn{7}{l}{\textit{*** $p < 0{,}001$; boundary test: $\frac{1}{2}\chi^2_0 + \frac{1}{2}\chi^2_1$.}}\\
\end{tabular}
\end{table}
